\begin{example}\label{ex:normalization-distance-projective}
The pseudometric retains only the mass in a closed K\"ahler background. Indeed \eqref{eq:normalization-closed-formula} gives
\[
 \Delta(\varphi,\psi)=|\log m_\varphi-\log m_\psi|.
\]
For an explicit calculation, let $\theta$ be the Fubini--Study form on $\mathbb P^n$, normalized by $\int_X\theta^n=1$, and let $D$ be a hyperplane. A defining section with its Fubini--Study metric gives $a_D=\log|s_D|^2\leq0$ after rescaling, with
\[
 \ddc a_D=[D]-\theta,\qquad
 \theta+\ddc(ta_D)=(1-t)\theta+t[D].
\]
Take a rational $0<t<1$ and put $\varphi_{t,D}=P_\theta[ta_D]$. As in \cref{ex:analytic-reference-normalization}, this is an admitted model reference of ordinary type $ta_D$. Off $D$ its original representative has Monge--Amp\`ere measure $(1-t)^n\theta^n$, and its non-pluripolar measure puts no mass on $D$. Invariance of mass under bounded changes in the closed background therefore gives
\[
 m_{\varphi_{t,D}}=(1-t)^n,\qquad
 B_{\varphi_{t,D}}(G)-B_0(G)=n\log(1-t).
\]
Thus the constant in \cref{prop:reference-normalization-comparison} is sharp.

On $\mathbb P^1$, choose two distinct points $p,q$ for $D$. The corresponding references have distance zero, although one is unbounded at $p$ and bounded at $q$, while the other has the reverse behavior. They are not ordinarily equivalent. For every positive integer $m$,
\[
 \mathcal I(m\varphi_{t,p})=\mathcal O(-\lfloor mt\rfloor p),\qquad
 \mathcal I(m\varphi_{t,q})=\mathcal O(-\lfloor mt\rfloor q).
\]
The one-variable integrability criterion proves these identities. For $mt\geq1$ the ideals differ. Thus $\Delta$ does not separate ideal types either, and cannot replace the singularity topology sought in this chapter.
\end{example}
